\chapter{Architettura del software}
\label{cap:architettura}

\section{Moduli}

Il codice \`e diviso per responsabilit\`a, non per file: ogni cartella di
\file{src/} ha un compito unico e dipende solo da quelle sotto di s\'e.

\begin{table}[H]
\centering
\footnotesize
\caption{Struttura di \file{src/}.}
\label{tab:moduli}
\begin{tabular}{@{}p{0.20\linewidth}p{0.72\linewidth}@{}}
\toprule
\textbf{Modulo} & \textbf{Responsabilit\`a} \\
\midrule
\file{common/} & tipo scalare unico (\code{scalar\_t}), tolleranza di
  validazione, utility di terminazione e cronometro \\
\file{index/} & partizionamento a blocchi degli indici globali: dimensione,
  offset e proprietario di ogni blocco \\
\file{gen/} & generatore riproducibile senza stato (\cref{sec:generazione}) \\
\file{serial/} & implementazione seriale minimale, usata come oracolo \\
\file{kernel/} & \code{local\_gemm}: un file per implementazione, \code{.c} o
  \code{.cu}, tutte dietro la stessa interfaccia \\
\file{mpi/} & griglia cartesiana, layout dei blocchi, distribuzione dei dati,
  algoritmo distribuito $Y = AX$ \\
\file{bench/} & driver di misura e validazione, parsing delle opzioni, output CSV \\
\bottomrule
\end{tabular}
\end{table}

La dipendenza pi\`u importante \`e quella che \textbf{non} esiste: \file{mpi/} non
include mai un header di un backend specifico. Vede solo
\file{kernel/kernel.h}, e non ha modo di sapere se sotto giri lo schema~A
scalare o un kernel CUDA.

\section{L'interfaccia unica del kernel locale}
\label{sec:interfaccia}

Ogni backend implementa tre funzioni di esecuzione e un insieme di canali di
misura e di introspezione:

\begin{lstlisting}[language=C,caption={Interfaccia di \file{src/kernel/kernel.h}.}]
/* esecuzione */
local_gemm_t *local_gemm_create(int m, int n, int k,
                                const scalar_t *A_loc,
                                int lda, int ldx, int ldy);
void local_gemm(local_gemm_t *ctx,
                const scalar_t *RESTRICT X, int ldx,
                scalar_t *RESTRICT Y, int ldy);
void local_gemm_destroy(local_gemm_t *ctx);

/* tempi della singola invocazione */
double local_gemm_last_compute_seconds(const local_gemm_t *ctx);
double local_gemm_last_h2d_X_seconds(const local_gemm_t *ctx);
double local_gemm_last_d2h_Y_seconds(const local_gemm_t *ctx);

/* byte trasferiti, per trasformare i tempi in banda */
size_t local_gemm_bytes_h2d_A(const local_gemm_t *ctx);
size_t local_gemm_bytes_h2d_X_per_call(const local_gemm_t *ctx);
size_t local_gemm_bytes_d2h_Y_per_call(const local_gemm_t *ctx);

/* preparazione, totale e scomposta */
double local_gemm_setup_seconds(const local_gemm_t *ctx);
double local_gemm_setup_device_init_seconds(const local_gemm_t *ctx);
double local_gemm_setup_device_alloc_seconds(const local_gemm_t *ctx);
double local_gemm_setup_h2d_A_seconds(const local_gemm_t *ctx);

/* piano di esecuzione effettivamente scelto a runtime */
int local_gemm_blocks_per_sm(const local_gemm_t *ctx);
int local_gemm_x_rows_per_tile(const local_gemm_t *ctx);

const char *kernel_name(void);
\end{lstlisting}

I canali oltre i primi tre sono la parte che ha richiesto pi\`u lavoro, e la
ragione \`e in \cref{sec:metriche}: la traccia esclude dal tempo $T$ i
trasferimenti da e verso la scheda e il preprocessing, ma consente
esplicitamente di misurarli e discuterli a parte --- e per discuterli non basta
il totale, servono le singole voci. Tutti seguono la stessa convenzione: un
valore \textbf{negativo} (o $0$ per i byte) significa ``non applicabile a questo
backend'', ed \`e ci\`o che rende una riga di CPU riconoscibile a colpo d'occhio
in un CSV che contiene anche righe GPU.

Le ultime due funzioni non sono metriche di prestazione ma la
\emph{configurazione} scelta a runtime dal backend: senza di esse una campagna
sul tiling mostrerebbe curve che salgono e scendono senza poter dire se sia
cambiato il tile, l'occupancy, o entrambi (\cref{sec:piano-tile}).

La scelta del backend \`e del \code{Makefile} (variabile \code{KERNEL}), non del
codice: non esiste nessun \code{switch} a runtime fra implementazioni, e quindi
nessun costo n\'e nessuna possibilit\`a di misurare per sbaglio un percorso
diverso da quello che si crede.

\subsection{Perch\'e il ciclo di vita ha tre fasi}
\label{sec:tre-fasi}

La separazione fra \code{create}, \code{local\_gemm} e \code{destroy} \`e la
scelta architetturale con l'effetto pi\`u diretto sulle misure, e discende da una
asimmetria del problema:

\begin{center}
\emph{$A$ non cambia mai fra un'invocazione e l'altra; $X$ e $Y$ cambiano sempre.}
\end{center}

$X$ arriva da un \code{MPI\_Bcast} e $Y$ deve tornare a un \code{MPI\_Reduce},
quindi entrambe attraversano ogni invocazione. $A$, invece, viene generata (o
ricevuta) una volta sola e resta identica per tutte le repetition. Tutto ci\`o che
riguarda $A$ pu\`o dunque stare in \code{create}, che \`e preprocessing e non viene
cronometrato.

Per il backend di CPU questo \`e quasi vuoto: \code{create} registra forma e
puntatore. Per i backend CUDA \`e invece il punto in cui $A$ viene copiata in VRAM
\emph{una volta sola}. Il guadagno non \`e cosmetico: con $A$ dell'ordine dei
gigabyte, una \code{cudaMemcpy} H2D per invocazione misurerebbe la banda PCIe e
non la GPU. La traccia consente esplicitamente di escludere dal tempo $T$ i
trasferimenti da e verso la scheda, e l'architettura rende questa esclusione una
propriet\`a strutturale invece che una sottrazione a posteriori.

Nella stessa fase vengono allocati anche i buffer di $X$ e $Y$ sul device e viene
eseguita una \code{cudaDeviceSynchronize} finale. La conseguenza \`e che nessuna
allocazione, e nessuna coda residua della copia di $A$, pu\`o cadere dentro la
prima repetition, neppure con \code{--warmup 0}.

\subsection{Lo stato in un contesto opaco}

Lo stato del backend vive in una \code{struct local\_gemm\_context} opaca, non in
variabili statiche di modulo. Il tipo \`e dichiarato in \file{kernel.h} e definito
in ciascun \code{.c}/\code{.cu}, quindi ogni backend gli d\`a i campi che gli
servono: il backend di CPU tiene forma e puntatore ad $A$, i backend CUDA
tengono i tre puntatori device, i due \code{cudaEvent} e i tempi.

Tre conseguenze pratiche: \`e esplicito nel tipo chi possiede la copia di $A$;
non esiste inizializzazione globale nascosta; il backend resta riutilizzabile,
perch\'e nulla vieta di avere pi\`u contesti vivi contemporaneamente.

\subsection{Compatibilit\`a C / C++ e il test che la difende}

I file \code{.cu} sono compilati da \code{nvcc} come C++, mentre il resto del
progetto \`e C11 compilato da \code{mpicc}. Due differenze toccano proprio
l'header condiviso: \code{restrict} non \`e una parola chiave del C++
(si scrive \code{\_\_restrict\_\_}) e i nomi delle funzioni C++ vengono decorati,
quindi senza \code{extern "C"} un backend \code{.cu} compilerebbe ma non
linkerebbe.

Entrambe le protezioni sono nell'header. La parte interessante \`e che la
propriet\`a \`e \emph{verificata}, e in due modi: il target \code{make check-cxx}
compila una sonda C++ che include \file{kernel.h} e \file{util.h}, e poi
controlla con \code{nm} che ognuno dei simboli dell'interfaccia --- tutte e
sedici le funzioni del listato precedente, pi\`u \code{xmalloc} e \code{die} ---
compaia \emph{non decorato} fra gli indefiniti dell'oggetto prodotto. Un
\code{extern "C"} dimenticato su una funzione aggiunta in seguito fallisce
quindi qui, su qualunque macchina e anche senza \code{nvcc}, e non come errore
di link sul server il giorno della campagna.

\section{Il driver di misura}

\file{src/bench/main.c} \`e l'unico punto che conosce sia MPI sia il kernel sia
le metriche. Fa, nell'ordine:

\begin{enumerate}[nosep,left=0pt]
  \item parsing e validazione delle opzioni (\cref{app:cli});
  \item creazione della griglia cartesiana e calcolo del layout locale;
  \item generazione o distribuzione di $A$ e $X$, e conversione del layout di
        $X$ se la build lo richiede --- ogni fase cronometrata separatamente
        (\cref{sec:tempi-esclusi});
  \item \code{local\_gemm\_create} (preprocessing del backend);
  \item \code{--warmup} invocazioni non cronometrate;
  \item \code{--reps} invocazioni cronometrate, con i tempi delle tre fasi
        raccolti per ogni repetition;
  \item riduzione dei tempi fra i rank, calcolo di media, mediana, minimo,
        deviazione standard e coefficiente di variazione;
  \item validazione opzionale contro l'oracolo seriale (\code{--check});
  \item stampa leggibile oppure riga CSV.
\end{enumerate}

Il fatto che warm-up e misura siano due cicli distinti sullo stesso codice, e non
un ciclo unico con i primi giri scartati a posteriori, \`e voluto: le
repetition cronometrate non contengono nessun ramo che le distingua dalle altre.

\section{Build: una configurazione, un binario}
\label{sec:build}

Ogni scostamento dalla configurazione di default produce un binario con un nome
proprio, come mostrato in \cref{tab:binari}. Non esiste un solo
\file{bin/matmul\_mpi} riscritto a ogni \code{make}.

\begin{table}[H]
\centering
\footnotesize
\caption{Corrispondenza fra configurazione di build e binario prodotto.}
\label{tab:binari}
\begin{tabular}{@{}p{0.46\linewidth}p{0.46\linewidth}@{}}
\toprule
\textbf{Comando} & \textbf{Binario} \\
\midrule
\code{make} & \file{bin/matmul\_mpi} \\
\code{make PREC=float} & \file{bin/matmul\_mpi-float} \\
\code{make FORCE\_GENERIC\_K=1} & \file{bin/matmul\_mpi-generic} \\
\code{make TEST\_A\_PADDING=8} & \file{bin/matmul\_mpi-pad8} \\
\code{make KERNEL=cuda\_naive} & \file{bin/matmul\_mpi-cuda\_naive} \\
\code{make KERNEL=cuda\_warp} & \file{bin/matmul\_mpi-cuda\_warp} \\
\code{make KERNEL=cuda\_warp\_smem} & \file{bin/matmul\_mpi-cuda\_warp\_smem} \\
\code{make KERNEL=cuda\_warp\_smem SMEM\_PAD=0} & \file{bin/matmul\_mpi-cuda\_warp\_smem-smempad0} \\
\code{make KERNEL=cuda\_warp\_smem TILE\_GRANULARITY=8} & \file{bin/matmul\_mpi-cuda\_warp\_smem-g8} \\
\code{make KERNEL=cuda\_warp\_smem SMEM\_BUDGET\_BYTES=16384} & \file{bin/matmul\_mpi-cuda\_warp\_smem-bud16384} \\
\code{make KERNEL=cuda\_warp\_tiled WARP\_COL\_TILE=4} & \file{bin/matmul\_mpi-cuda\_warp\_tiled-tile4} \\
\code{make KERNEL=cuda\_warp X\_LAYOUT=column} & \file{bin/matmul\_mpi-cuda\_warp-xcol} \\
\code{make BLOCK=512 KERNEL=cuda\_warp} & \file{bin/matmul\_mpi-cuda\_warp-blk512} \\
\code{make KERNEL=cublas} & \file{bin/matmul\_mpi-cublas} \\
\bottomrule
\end{tabular}
\end{table}

I parametri che entrano nel nome sono solo quelli che il backend selezionato
legge davvero: \code{TILE\_GRANULARITY} e \code{SMEM\_BUDGET\_BYTES} compaiono
soltanto con \kernelname{cuda\_warp\_smem}, \code{WARP\_COL\_TILE} soltanto con
\kernelname{cuda\_warp\_tiled}. Altrimenti si otterrebbero binari con nomi
diversi e contenuto identico, indistinguibili nel CSV perch\'e
\code{kernel\_name()} non cambierebbe. Anche gli oggetti sono per
configurazione (\file{obj/matmul\_mpi<config>/}), quindi cambiare precisione o
backend non pu\`o riutilizzare per sbaglio una compilazione precedente.

Il motivo \`e sperimentale, non organizzativo: una build non pu\`o sovrascriverne
un'altra e far misurare in singola precisione una campagna che si credeva in
doppia. Nello stesso spirito, \code{kernel\_name()} riporta nel CSV anche i
parametri che cambiano il codice generato --- per esempio
\code{cuda\_warp\_tiled(tile8)} oppure il suffisso \code{(blk512)} quando la
dimensione del blocco non \`e quella di default --- cos\`i le righe di uno sweep
restano distinguibili anche a distanza di settimane.

Il \code{Makefile} riconosce da s\'e se il backend richiesto \`e C o CUDA: se
esiste \file{src/kernel/<nome>.cu} lo compila con \code{nvcc -arch=sm\_75} e
linka \code{-lcudart} (e \code{-lcublas} per il backend cuBLAS). Su una macchina
priva di \code{nvcc} la build CUDA si ferma con un messaggio esplicito, mentre
tutto il resto del progetto continua a compilare: \`e la propriet\`a che ha
permesso di sviluppare su una macchina ARM64 senza GPU e misurare sul server.
